\section{Time-Varying Coupled Analysis}

\subsection{Introduction}
This chapter details the physics of time-varying engine performance, focusing on the coupled evolution of geometry, thermodynamics, and stability. Unlike static analysis, time-varying analysis accounts for the cumulative effects of ablation, propellant depletion, and transient operating conditions over the full burn duration.

\subsection{Key Files Referenced}
The physics models documented in this chapter are implemented in:
\begin{itemize}
    \item \texttt{engine/core/runner.py} - Time-series evaluation (\texttt{evaluate\_arrays\_with\_time})
    \item \texttt{engine/pipeline/thermal/ablative\_geometry.py} - Geometry evolution from material loss
    \item \texttt{engine/pipeline/localized\_ablation.py} - Enhanced ablation at impingement zones
    \item \texttt{engine/pipeline/time\_varying\_solver.py} - Fully-coupled transient solver
\end{itemize}

\subsection{Nomenclature}
\begin{description}
    \item[$\dot{r}$] Recession rate [\si{m/s}]
    \item[$H_{abl}$] Heat of ablation [\si{J/kg}]
    \item[$V_c(t)$] Time-varying chamber volume [\si{m^3}]
    \item[$A_t(t)$] Time-varying throat area [\si{m^2}]
    \item[$\epsilon(t)$] Time-varying expansion ratio
    \item[$q''_{imp}$] Impingement heat flux [\si{W/m^2}]
\end{description}

\subsection{Geometry Evolution}
The engine geometry is updated at each time step based on the cumulative material loss from the ablative liner and graphite insert.

\subsubsection{Chamber Volume and Throat Area}
The instantaneous chamber volume and throat area are calculated by integrating the local recession rates:
\begin{equation}
D(x, t) = D_0(x) + 2 \int_0^t \dot{r}(x, \tau) d\tau
\end{equation}
\begin{equation}
A_t(t) = \frac{\pi}{4} D_t(t)^2, \quad V_c(t) = \int_0^{L_c} \frac{\pi}{4} D(x, t)^2 dx
\end{equation}

\subsection{Localized Ablation (Fuel Impingement)}
For pintle injectors, the fuel spray creates localized zones of enhanced heat flux where it impinges on the chamber wall.

\subsubsection{Impingement Zone Modeling}
The impingement heat flux is modeled using a Gaussian distribution centered at the impingement point $L_{imp}$:
\begin{equation}
q''_{total}(x) = q''_{base}(x) \cdot \left[1 + (M_{imp} - 1) \cdot \exp\left(-0.5 \left(\frac{x - L_{imp}}{\sigma_{imp}}\right)^2\right)\right]
\end{equation}
where $M_{imp}$ is the peak multiplier (typically 2.0--5.0x).

\subsection{Coupled Solver Logic}
The time-varying solver iterates through the burn duration, performing the following steps at each interval $\Delta t$:
\begin{enumerate}
    \item \textbf{Static Evaluation}: Solve for $P_c, \dot{m}, F$ using current geometry $V_c(t), A_t(t)$.
    \item \textbf{Thermal Analysis}: Compute local heat fluxes $q''(x)$ and recession rates $\dot{r}(x)$.
    \item \textbf{Stability Check}: Evaluate stability margins for the current geometry and operating point.
    \item \textbf{State Update}: Update cumulative recession and advance geometry for $t + \Delta t$.
\end{enumerate}

\subsection{Transient Effects on Performance}
As the throat area $A_t$ increases due to recession, the chamber pressure $P_c$ typically decreases for a fixed propellant supply. This leads to:
\begin{itemize}
    \item \textbf{Thrust Decay}: Reduction in momentum thrust.
    \item \textbf{Expansion Loss}: Shift in expansion ratio $\epsilon = A_e/A_t$, leading to potential over-expansion or under-expansion losses.
    \item \textbf{Stability Shifts}: Changes in acoustic frequencies and chugging margins as $L^*$ and $V_c$ evolve.
\end{itemize}
